% XeTeX's \XeTeXlinebreaklocale: ICU's line breaks for the locale (its
% rule sets: root's, Japanese and Korean's, Chinese's, the lb= variants),
% each a \XeTeXlinebreakpenalty and \XeTeXlinebreakskip (glue only if not
% zero_glue, the penalty if not zero or no glue); words in a native font
% (Latin Modern: CJK glyphs missing, as in the oracle; breaks do not
% depend on glyphs).
\catcode`\{=1 \catcode`\}=2 \catcode`\#=6 \catcode`\^=7
\showboxdepth=100 \showboxbreadth=100000 \tracingonline=1 \tracinglostchars=0
\hsize=100pt \vsize=600pt \parfillskip=0pt plus 1fil \baselineskip=12pt
\tolerance=10000 \parindent=0pt \output={\shipout\box255}
\font\x="[lmroman10-regular.otf]" \x
\def\t#1#2{\XeTeXlinebreaklocale "#1"\relax #2\showlists\par\showlists}
\t{ja}{「こんにちは」、世界。ちょっと待ってください……ねえ、ゃゅょっァィ}
\t{zh}{这是一个测试：中文、英文（English）和数字123混排。《三国演义》}
\t{zh_TW}{這是一個測試：中文、英文（English）和數字123混排。}
\t{ko}{한국어 문장입니다. 띄어쓰기가있습니다!}
\t{en}{http://example.com/path?x=1&y=2 e.g. U.S.A. AT&T don't 10% $10.00-5}
\t{th}{東京は English 123}
\t{}{「こんにちは」、世界。}
\t{ja@lb=loose}{ラーメン・カレー・ソバ〜！？々ゝゞヽヾーァ}
\t{ja@lb=strict}{ラーメン・カレー・ソバ〜！？々ゝゞヽヾーァ}
\t{zh-u-lb-normal}{价格是NT$1,234.56元，约US$40。ー〜}
\t{en@lb=loose}{ラーメン・カレー（123）}
\t{ja}{東京は２０２６年１０月６日（火）です。😀👍🏽𠀋𡈽}
\XeTeXlinebreakpenalty=50
\t{ja}{日本語の文章です。句読点、括弧（かっこ）}
\XeTeXlinebreakskip=0pt plus 1pt
\t{ja}{日本語の文章です。句読点、括弧（かっこ）}
\XeTeXlinebreakpenalty=0
\t{zh}{这是一个测试：中文、英文}
\XeTeXlinebreakskip=0pt
\t{zh}{这是一个测试：中文、英文}
\XeTeXlinebreakskip=0pt plus 0fil
\t{ja}{東京は２０２６年}
% a name goes on over any character, to a space
\setbox0\hbox{\XeTeXlinebreaklocale "ja"東京は２０２６年} \showbox0
\XeTeXlinebreaklocale "ja" a\-b-c東京\-は--大阪 東 東京\par
\begingroup \XeTeXlinebreaklocale "zh" \endgroup 中文测试。\par
\setbox0\vbox{\XeTeXlinebreaklocale "ja" \XeTeXlinebreakskip=0pt plus 2pt
  吾輩は猫である。名前はまだ無い。どこで生れたかとんと見当がつかぬ。何でも薄暗いじめじめした所でニャーニャー泣いていた事だけは記憶している。}
\showbox0 \box0
\XeTeXlinebreaklocale "" \XeTeXlinebreakskip=0pt
\end
